
\begin{frame}{Proposta TCC: Modelagem de Redes Complexas no Esporte Olímpico}

    \vspace{0.5cm}
    \begin{block}{Área de Pesquisa}
        \textbf{Análise de Dados} \\
        \textit{(De acordo com a tabela CAPES/CNPQ: 1.02.02.08-0)}
    \end{block}
    \vspace{0.5cm}
    
\end{frame}

\begin{frame}{Tema}
    \begin{block}{Tema}
            \textbf{Aplicação de redes complexas para modelar e analisar interações competitivas entre atletas olímpicos.}
            \vspace{0.3cm}
            \textit{Motivação: Revelar padrões e conexões entre atletas que análises tradicionais não conseguem identificar, oferecendo novos insights sobre desempenho e competição.}
            \begin{itemize}
                \item Utilização de dados históricos de medalhas, categorias esportivas e eventos.
                \item Identificação de padrões estruturais e dinâmicas de colaboração/competição.
                \item Métodos de redes complexas para análise de sistemas interconectados.
            \end{itemize}
        \end{block}
\end{frame}

\begin{frame}{Objetivos da Modelagem}
    \framesubtitle{Caracterização das Interações Competitivas}
    \begin{block}{Objetivo Geral}
        \textbf{Caracterizar as interações competitivas entre atletas olímpicos através de métricas de redes complexas.}
        \begin{itemize}
            \item Foco na identificação de padrões de conectividade.
            \item Identificação de fatores determinantes de sucesso.
            \item Análise de dinâmicas colaborativas.
        \end{itemize}
    \end{block}
    \vspace{0.3cm}
    
\end{frame}

\begin{frame}{Objetivos da Modelagem}
\begin{block}{Objetivos Secundários}
        \begin{itemize}
            \item \textbf{Modelagem de Grafos:} Direcionados e ponderados, representando conexões de pódios compartilhados.
            \item \textbf{Aplicação de Métricas:} Centralidade, modularidade e coeficiente de agrupamento.
            \item \textbf{Exploração de Diferenças Estruturais:} Entre esportes individuais e coletivos, análise de comunidades.
            \item \textbf{Visualizações Interativas:} Facilitação da compreensão dos resultados.
        \end{itemize}
    \end{block}
\end{frame}


\begin{frame}{Base de dados}
\begin{block}{Base de Dados}
    \begin{itemize}
        \item \textbf{Nome:} 120 years of Olympic history.
        \item \textbf{Descrição:} Basic bio data on athletes and medal results from Athens 1896 to Rio 2016.
        \item \textbf{Autor:} Randi H Griffin, PhD.
        \item \textbf{Disponível em:} \url{https://www.kaggle.com/datasets/heesoo37/120-years-of-olympic-history-athletes-and-results}.
    \end{itemize}
\end{block}
\end{frame}

\begin{frame}[fragile]{Caracteristicas da Base de Dados}
\framesubtitle{Premissa de Análise}
\begin{block}{Características da Base de Dados}
    \begin{itemize}
        \item \textbf{18 Colunas.}
        \item \textbf{271117 Linhas}
        \item \textbf{Campos:} ID, Name, Sex, Age, Height, Weight, Team, NOC, Games, Year, Season, City, Sport, Event, Medal.
    \end{itemize}
\end{block}
\end{frame}

